
\section*{AI Use Statement}

We used generative AI tools (Opus 4.8, Opus 5.0 and Fable 5) as coding and writing assistants throughout this work. Specifically: for implementing and refactoring the data pipeline, the grading harness, the fine-tuning scripts, and the analysis code; for drafting and editing prose in this paper; and for exploratory literature search. We did not use generative AI to generate research ideas, to produce experimental results, or to write any of the analysis numbers reported here --- every number in this paper is computed by the released scripts from the released data, and the ledgers those scripts emit are the authority for the manuscript. All AI-assisted code was reviewed and executed by the authors, and the statistical results were independently re-derived before being reported. We take responsibility for the final content of this work, including all text, claims, and artifacts.

Note that generative models are also the \emph{object} of study here rather than a tool: the LLM graders evaluated in Sections~\ref{sec:closed}--\ref{sec:finetune} produced the scores we analyse, and those outputs are released in full.

\section*{Reproducibility Statement}

Every number in the paper is regenerated from released artifacts by released code. The two anonymised exam datasets, the rubrics and reference solutions, one result workbook per configuration, and the analysis, grading, and fine-tuning pipelines \artifactsshipin. Section~\ref{sec:setup} and Appendix~\ref{app:details} specify the prompt, the serving stack, the decoding parameters, and the exclusion rules; Appendix~\ref{app:runs} tabulates all $171$ CV-exam configurations with their bootstrap confidence intervals; Appendix~\ref{app:ft} gives the fine-tuning hyperparameters and the seed-fixed split. Training uses HuggingFace with PEFT and every evaluated number is served by vLLM, Gemma-4's adapter included. One caveat is stated where it arises: evaluation is not bitwise reproducible because vLLM batching reorders reductions ($\sim 0.01$ MAE); Gemma-4's vLLM serving path was cross-validated against HuggingFace generation to within $\sim 0.1$ MAE (Appendix~\ref{app:ft}).

\section*{Ethics Statement}

The data are student exam submissions and grader marks from two university courses, released with the course instructors' permission. Both datasets are anonymised before release: student identifiers, names, and email addresses are removed from notebook code, markdown, outputs, and file metadata, and graders appear only as opaque identifiers. No demographic attributes are collected or released, and no student is identifiable from the released artifacts.

The application this paper studies carries real risk to the people being graded, which is why we report the failure mode rather than only the headline accuracy. An under-grading collapse of the kind documented in Section~\ref{sec:brittle} harms students who attempted the work in good faith, and it is invisible to single-prompt evaluation. We therefore recommend against deploying an LLM grader on the basis of one prompt's measured accuracy, and we report bias alongside MAE throughout so that the direction of harm is visible. We regard human oversight of consequential grading decisions as a requirement, not an option, and none of the results here should be read as licensing unsupervised automated assessment.
